\section*{Avant-propos}
\addcontentsline{toc}{section}{\protect\numberline{}Avant-propos}
\label{sec:avant_propos}

Le présent document constitue un document de suivi de projet complet relatif au développement du lanceur expérimental bi-étage
SP-02, mené au sein de l'association AéroIPSA et destiné à être présenté dans le cadre du programme \gls{Fusex} de
\gls{Planète Sciences}. Il a pour vocation d'assurer la traçabilité des choix techniques, de l'organisation du projet et de la
maîtrise des risques tout au long du cycle de conception.\\

Ce document remplit plusieurs fonctions complémentaires. D'une part, il constitue un support d'échange et d'évaluation à
destination de \gls{Planète Sciences} et de ses suiveurs de projet, notamment dans le cadre des revues techniques \acrfull{RCE},
pour démontrer la conformité du projet aux exigences du \acrfull{CDC} \gls{Fusex} bi-étage et la maîtrise de la sûreté de
fonctionnement.\\

D'autre part, le projet SP-02 sert de support au \acrfull{PMI} réalisé dans le cadre de la formation de la dernière année
d'ingénieur à l'\acrfull{IPSA}. Le \acrshort{PMI} constitue un projet académique encadré, borné dans le temps, répondant à une
problématique bien définie, et évalué selon des critères pédagogiques spécifiques. Dans le cas présent, le sujet de ce dernier
a été proposé par les membres en 5ème année inscrit au projet SP-02, porte la référence \og{}\acrfull{ET} numéro 2\fg{}, et est
intitulé :

\begin{quote}
\textit{\og{}Perspective d'expérience d'un système de séparation d'étages de fusée par aérofreinage\fg{}}.
\end{quote}

Le \acrshort{PMI} s'inscrit ainsi comme un sous-ensemble du projet SP-02, centré sur l'étude, la conception, la modélisation,
l'instrumentation et la validation expérimentale du système de séparation par aérofreinage, ainsi que des systèmes
électroniques de mesure et de commande associés. Les travaux relevant du \acrshort{PMI} font l'objet de livrables spécifiques
(fiches navettes, rapport final, soutenance) et suivent un calendrier académique distinct de celui du projet associatif global.\\

Ainsi, les différentes sections du document participent :
\begin{itemize}
    \item à la réponse aux exigences formelles du \acrshort{PMI} (démarche d'ingénierie, simulations \acrshort{CFD} et
    \acrshort{FEM}, livrables académiques, ...),
    \item au suivi global du projet SP-02 en tant que lanceur expérimental (conception des sous-systèmes, sûreté de
    fonctionnement, validation, essais, organisation, ...),
\end{itemize}
sans que ces deux dimensions ne soient dissociées.\\

Le choix a été fait de conserver un document unique, garantissant une vision
cohérente du projet, à la fois sur le plan technique, organisationnel et académique.\\

Le présent rapport s'adresse ainsi à plusieurs publics : les encadrants et jurys académiques du \acrshort{PMI}, les organisateurs
et suiveurs de \gls{Planète Sciences}, ainsi que les membres actuels et futurs de l'association AéroIPSA, chacun pouvant y trouver
les éléments correspondant à ses attentes et à son niveau d'analyse.

\newpage